\section{What the Entity Route Costs}
\label{sec:ch06-what-the-entity-route-costs}

\index{entity route!compared with the node route|(}%
Two routes now reach the same row. Ask for one session's title through
\code{node(id:)} and through \code{\_entities}, and compare what SQLite is sent.

\begin{minted}{text}
-- the node route, selecting title
SELECT "s"."Title"
FROM "Sessions" AS "s"
WHERE "s"."Id" = @id
LIMIT 1
\end{minted}

\begin{minted}[fontsize=\footnotesize]{text}
-- the entity route, selecting title
SELECT "s"."Id", "s"."Abstract", "s"."DurationMinutes", "s"."SpeakerId", "s"."StartsAt", "s"."Title"
FROM "Sessions" AS "s"
WHERE "s"."Id" = @key
LIMIT 1
\end{minted}

\index{projection!absent from the entity route}%
One column against six, for the same field of the same row. Nothing about the
query differed. What differed is the parameter list:
section~\ref{sec:ch06-the-key-arrives-as-a-string} gave the reference resolver
no \code{QueryContext}, so the projection
chapter~\ref{ch:data-without-n-plus-1} built does not reach it. Six columns
is the whole row of a small table. On a wide one it is every column you have,
fetched to answer a question about one of them.

\index{statement count!per representation}%
Then the counts, which chapter~\ref{ch:router-n-plus-1} spends a chapter on and
which are small enough here to read in one line each.

\begin{minted}{text}
3 Session representations  ->  3 statements
3 Speaker representations  ->  1 statement
\end{minted}

\index{DataLoader!behind a reference resolver}%
Same field, same call, same number of keys. The difference is one line in each
resolver: \code{SessionType} goes to the table and \code{SpeakerType} goes
through the DataLoader chapter~\ref{ch:data-without-n-plus-1} already had. A
loader behind a reference resolver batches a whole call into one statement,
because the package resolves the representations together and the loader
collects them, exactly as it does for a list of sibling fields.

\index{entities@\code{\_entities}!batches its representations}%
Which is the fact chapter~\ref{ch:router-n-plus-1} is built on, and it is
already true of a subgraph with nothing in front of it. The entity route is a
batching interface. A router uses it that way on purpose, collecting a whole
level of a query plan into one call rather than asking once per object, and
whether that batch costs one statement or one per representation is decided
inside the method you wrote before any router exists to send it. Both of this
chapter's are correct as written.

\index{figure!the two routes}%
Figure~\ref{fig:ch06-two-routes} is the pair of routes drawn against each other.

\begin{figure}
  \centering
  % Two routes to the same row: node(id:), built in chapter 4, and the
% _entities(representations:) reference resolver this chapter adds.
%
% Same idiom as figures/tikz/ch01-one-field.tex and
% figures/tikz/ch05-the-join-moves.tex: each lane is a timeline read
% left to right, ticks mark stages, and a brace under the stage that
% matters carries a one-line label. The four ticks sit at the same x
% in both lanes, because the argument is that the two routes do the
% same work at the same points and differ only in what they are
% handed and in where the decoding happens.
%
% The one asymmetry the figure draws: in lane A, decoding is its own
% stage, tick 2, run by Hot Chocolate before any resolver is entered.
% Lane B's tick 2 is the same point in the sequence and nothing happens
% there, which its label says; the id string arrives whole and the
% reference resolver, tick 3, decodes it itself. The brace therefore
% lands under tick 2 in lane A and tick 3 in lane B, and that shift is
% what the eye should catch first.
%
% Lane B's SELECT elides the middle four columns and states the count
% in words, keeping the first and last column named. Both lanes cost
% one statement for one key; neither is drawn as cheaper.
%
% No quotation marks anywhere in this file. SPEC decision 30 puts the
% book in strict Ascii mode and the prose gate reads a literal " as a
% quote character, while \" is LaTeX's diaeresis accent and fails to
% compile inside a node. The labels are written so that neither is
% needed: identifiers appear bare.
%
% Bare tikzpicture (house style): the figure environment, caption and
% label live at the call site in the section file.
\begin{tikzpicture}[
  x=1mm, y=1mm,
  every node/.append style={font=\sffamily\scriptsize},
  axis/.style={draw, line width=0.4pt,
    -{Straight Barb[length=1.4mm, width=1.6mm]}},
  tick/.style={draw, line width=0.4pt},
  event/.style={anchor=south, align=center, text width=24mm,
    inner sep=1pt},
  span/.style={draw, line width=0.4pt},
  spanlabel/.style={anchor=north, align=center, inner sep=2pt},
  lane/.style={anchor=west, font=\sffamily\scriptsize\itshape},
]

% ------------------------------------------------------------- lane A
\node[lane] at (0,20)
  {A. the node route, asking for one field};

\draw[axis] (0,0) -- (150,0);
\foreach \x in {16,58,100,138}{
  \draw[tick] (\x,-1.5) -- (\x,1.5);
}

\node[event, text width=30mm] at (16,3)
  {client sends\\node(id: U2Vzc2lvbjox),\\selecting title};
\node[event, text width=28mm] at (58,3)
  {the id is decoded\\before the resolver\\is entered};
\node[event, text width=30mm] at (100,3)
  {node resolver is given\\the integer 1 and a\\QueryContext};
\node[event, text width=26mm] at (138,3)
  {SELECT s.Title:\\one column};

\draw[span] (44,-4) -- (44,-7) -- (72,-7) -- (72,-4);
\node[spanlabel] at (58,-8)
  {decoded here, by the framework,\\as a stage of its own};

% ------------------------------------------------------------- lane B
\begin{scope}[shift={(0,-46)}]
  \node[lane] at (0,20)
    {B. the entity route, asking for the same field};

  \draw[axis] (0,0) -- (150,0);
  \foreach \x in {16,58,100,138}{
    \draw[tick] (\x,-1.5) -- (\x,1.5);
  }

  \node[event, text width=32mm] at (16,3)
    {router sends one\\representation:\\\_\_typename Session,\\id
     U2Vzc2lvbjox};
  \node[event, text width=26mm] at (58,3)
    {nothing decodes\\anything};
  \node[event, text width=30mm] at (100,3)
    {reference resolver is\\given the string, and\\no QueryContext};
  \node[event, text width=26mm] at (138,3)
    {SELECT s.Id, ...,\\s.Title:\\six columns};

  \draw[span] (86,-4) -- (86,-7) -- (114,-7) -- (114,-4);
  \node[spanlabel] at (100,-8)
    {decoded here, by hand, inside\\the resolver you wrote};
\end{scope}

\end{tikzpicture}

  \caption{The same row, reached twice. The stages line up because the two
  routes do the same work in the same order; what differs is what each resolver
  is handed. Lane~A is given a decoded integer and a projection, and asks for
  one column. Lane~B is given a string, decodes it itself, is given no
  projection, and asks for six. The braces mark where the decoding happens, and
  they sit under different stages on purpose: that shift is the only structural
  difference between the lanes, and the source of both others.}
  \label{fig:ch06-two-routes}
\end{figure}

\index{entity route!is not the node route}%
It is tempting to read the entity route as \code{node(id:)} with a different
spelling, since both take a global id and return the object it names. They are
different code with different properties, and this chapter has now found three
differences: what the resolver is handed, what it takes, and what it costs.
Chapter~\ref{ch:entities-authorization} finds the fourth and most expensive one.
\index{entity route!compared with the node route|)}%
